\begin{problem}[33届物理竞赛决赛]{68}
    五、许多赛车上都装有可调节的导流翼片，可以为水平道路上的赛车提供竖直向上或向下的附加压力。如果赛车速度的大小为 $v$，则上述压力的大小为 $f_{B} = c_{B} v^{2}$ ，$c_{B}$ 为一常量。当导流翼片的前方上翘时，压力方向向上；当导流翼片的后方上翘时，压力方向向下。赛车在运动过程中受到迎面空气的阻力，阻力大小为 $f_{A} = c_{A} v^{2}$ ，$c_{A}$ 为一常量。已知赛车质量为 $m$，轮胎与路面之间的静摩擦系数为 $\mu_{S}$ （$\mu_{S}<1$ ）。
    \begin{subproblem}
        赛车在水平直道上匀速行驶时，考虑到在运动过程中轮胎的形变，路面对赛车会形成阻力，阻力大小与车对路面的正压力大小成正比，比例系数为 $\mu_{R}$ （$\mu_{R} < \mu_{S}$ ）。若导流翼片被调至前方上翘，求当赛车行驶速度大小为多大时，赛车发动机输出的功率最大？最大输出功率为多少？
    \end{subproblem}
    \begin{subproblem}
        不考虑赛车在运动过程中轮胎的形变所引起的地面阻力，求当赛车在水平地面内沿半径为 $r$ 的圆形道路上匀速率行驶、且不沿路面发生滑动或飞离地面时所允许的速率最大值 $v_{\max}$ 。在这种情况下，导流翼片应被调至前方上翘还是后方上翘，$v_{\max}$ 可以更大？假设赛车发动机输出的功率可以足够大。
    \end{subproblem}
\end{problem}